\documentclass[11pt]{article}
\usepackage[a4paper,margin=1.85cm]{geometry}
\usepackage[T1]{fontenc}
\usepackage{lmodern}
\usepackage{amsmath,amssymb,amsthm,mathtools}
\usepackage{graphicx}
\usepackage{xcolor}
\usepackage{microtype}
\usepackage{enumitem}
\usepackage[hidelinks]{hyperref}

\definecolor{ink}{HTML}{172033}
\definecolor{accent}{HTML}{7A3E38}
\definecolor{soft}{HTML}{FAF0EE}
\color{ink}
\newtheorem{theorem}{Theorem}
\newtheorem{proposition}{Proposition}
\newtheorem{corollary}{Corollary}

\newcommand{\de}{\delta_e}

\title{A six-point counterexample to the idempotent-measure projectivity conjecture}
\author{Solution packet for the conjecture in arXiv:2203.09765}
\date{11 August 2026}

\begin{document}
\maketitle

\begin{center}
\fcolorbox{accent}{soft}{\parbox{0.91\textwidth}{
\textbf{Full negative answer.}
The conjecture fails already for the finite group $S_3$.  An idempotent measure
$\mu$ of norm $2$ generates exactly the same left ideal of $\ell^1(S_3)$ as a
norm-one idempotent $q$.  The ideal is therefore metrically projective, although
the unique density of $\mu$ has norm $2$, not $1$.}}
\end{center}

\section*{The conjecture}
After Proposition 3.2.4, Nemesh conjectures that for an idempotent measure
$\mu\in M(G)$ the left ideal $I=L^1(G)*\mu$ is metrically projective if and
only if $\mu=p\,m_G$ for some $p\in I$ with $\|p\|_1=1$ \cite{Source}.

\begin{center}
\includegraphics[width=0.94\textwidth]{figures/source_conjecture_crop.png}
\end{center}

\section*{Idea of the counterexample}
Metric projectivity belongs to the module $I$, whereas the proposed norm
condition belongs to the particular idempotent measure used to present $I$.
These data are not invariant together.  If $q$ is an idempotent and
$x\in(1-q)Aq$, then
\[
 qx=0,\qquad xq=x,\qquad x^2=0.
\]
Consequently $\mu=q+x$ is another idempotent, and
$A\mu=Aq$.  In a noncommutative group algebra one can choose $q$ with norm
one while arranging that $q$ and $x$ have disjoint supports, forcing
$\|\mu\|_1>1$.

\section*{The explicit example}
Let $G=S_3$, with products composed from right to left, and identify
$M(G)=L^1(G)=\ell^1(G)$ using counting Haar measure.  Put
\[
 s=(12),\qquad a=\delta_{(23)},\qquad
 q=\frac12(\de+\delta_s),\qquad
 x=(\de-q)*a*q,
\]
and set $\mu=q+x$.  Direct expansion gives
\begin{equation}\label{eq:x}
 x=\frac14\bigl(
 \delta_{(23)}+\delta_{(132)}-\delta_{(123)}-\delta_{(13)}
 \bigr).
\end{equation}

\begin{theorem}\label{thm:counterexample}
The measure $\mu$ is idempotent, the left ideal
$I=\ell^1(S_3)*\mu$ is a metrically projective $\ell^1(S_3)$-module, and
$\mu$ has no $L^1$-density of norm one.  Hence the conjectured ``only if''
direction is false.
\end{theorem}

\begin{proof}
Since $q$ is the normalized Haar measure of the subgroup
$\{e,(12)\}$, it is an idempotent and $\|q\|_1=1$.  The definition of $x$
immediately yields
\[
 q*x=0,\qquad x*q=x,
\]
and
\[
 x*x=(\de-q)*a*q*(\de-q)*a*q=0
\]
because $q*(\de-q)=0$.  Therefore
\[
 \mu*\mu=q*q+q*x+x*q+x*x=q+x=\mu.
\]
The same identities give
\begin{equation}\label{eq:relations}
 \mu*q=\mu,\qquad q*\mu=q.
\end{equation}
The first relation in \eqref{eq:relations} implies
$\ell^1(G)\mu\subseteq\ell^1(G)q$, and the second gives the reverse
inclusion.  Thus
\[
 I=\ell^1(G)\mu=\ell^1(G)q.
\]
Since $q$ is a norm-one idempotent, Proposition 2.1.13(i) of the source says
that $\ell^1(G)q$ is metrically projective.  Equivalently, it is the
contractive retract of the free rank-one left module $\ell^1(G)$ given by
right convolution with $q$ and the inclusion of its range.

The supports of $q$ and $x$ in \eqref{eq:x} are disjoint, so
\[
 \|\mu\|_1=\|q\|_1+\|x\|_1=1+1=2.
\]
For any choice of Haar normalization, the $L^1$ norm of the Radon--Nikodym
density of $\mu$ is its total-variation norm, namely $2$.  That density is
unique up to Haar-null sets.  Hence there is no $p$ with
$\mu=p\,m_G$ and $\|p\|_1=1$, even though $I$ is metrically projective.
\end{proof}

\section*{A general family}
The phenomenon is not special to the six-element computation.

\begin{proposition}\label{prop:family}
Let $G$ be a discrete group with a finite nonnormal subgroup $H$.  Write $q$
for normalized Haar measure on $H$, and choose
$g\notin N_G(H)$.  Then
\[
 x=(\de-q)*\delta_g*q,\qquad \mu=q+x
\]
give an idempotent $\mu\in\ell^1(G)$ such that
\[
 \ell^1(G)\mu=\ell^1(G)q,\qquad \|q\|_1=1,\qquad \|\mu\|_1>1.
\]
Thus every such group supplies a counterexample to the conjecture.
\end{proposition}

\begin{proof}
The Peirce relations used in Theorem~\ref{thm:counterexample} are formal and
give idempotency and equality of the two left ideals.  If $x=0$, then
$q*\delta_g*q=\delta_g*q$, so the right coset $gH$ is invariant under left
multiplication by $H$.  Hence $g^{-1}Hg\subseteq H$, and finiteness gives
$g^{-1}Hg=H$, contrary to $g\notin N_G(H)$.  Thus $x\ne0$.  Its support is
contained in the double coset $HgH$, which is disjoint from $H$; therefore
$\|\mu\|_1=\|q\|_1+\|x\|_1>1$.
\end{proof}

\section*{Adversarial checks}
\begin{enumerate}[leftmargin=2em]
\item The inclusions of principal \emph{left} ideals have the correct
orientation: $\mu q=\mu$ gives $A\mu\subseteq Aq$, while $q\mu=q$ gives
$Aq\subseteq A\mu$.
\item The source's sufficient criterion is applied to the common ideal using
the norm-one generator $q$, not to the norm-two generator $\mu$.
\item The density clause cannot be rescued by choosing another density:
Radon--Nikodym densities with respect to a fixed Haar measure are unique, and
their $L^1$ norm equals total variation.
\item All coefficients and convolution identities in the $S_3$ example were
checked with exact rational arithmetic; no floating-point or convention-
sensitive inference is used.
\end{enumerate}

\section*{Bounded novelty check}
On 11 August 2026, the arXiv id, exact conjecture, idempotent measures,
principal left ideals, and metric projectivity were checked in the four run
indexes and by bounded exact-phrase web/arXiv searches.  The conjecture also
appears unchanged in the recent published version located by the search.  No
later resolution or this $S_3$ counterexample was found.  This is evidence for
novelty, not a definitive guarantee.

\begin{thebibliography}{9}
\bibitem{Source}
N. Nemesh,
\emph{Metric and topological homology. Projectivity, injectivity and flatness},
arXiv:2203.09765, 2022.
\end{thebibliography}

\end{document}
